%!TEX root = ../hutbthesis_main.tex
\chapter{研究理论基础与关键技术}
\section{人体上肢生物力学理论}
人体上肢运动，是肌肉、骨骼和关节一起配合的复杂动力学过程，也是肌肉驱动交互系统最核心的生理基础；上肢以肩、肘、腕这几个关节为运动枢纽，以骨骼肌作为动力来源，通过肌肉收缩产生拉力，带动骨骼绕着关节轴线，完成转动、屈伸、旋转等动作，最后实现手部的精准操控；在生物力学模型里，一般会把上肢简化成多刚体的连杆结构，用 Hill 肌肉模型来描述肌肉收缩的特点；如图~\ref{img-1}的肌肉模型示意图所示，肌肉力可以表示为激活度、肌肉长度、收缩速度和肌腱弹性的综合函数；肌肉激活度是主要的控制输入，它决定了肌肉输出力的大小；不同肌群的协同激活和拮抗配合，能够实现稳定、柔顺又精准的末端运动；对于方向盘操控这类旋转类任务来说，上肢肌肉的对称发力、力矩平衡和连续力输出，是特别关键的；所以，建立符合生理特性的肌肉 — 关节动力学模型，是实现精准、自然人机交互的前提\cite{wang2026simulating}
\begin{figure}[htbp]
  \centering
  \includegraphics[width=0.8\textwidth]{image_2.png}
  \caption{Hill肌肉模型示意图}
  \label{img-1}
\end{figure}  





为定量描述上肢肌肉的力学输出特性，本研究采用 Hill 三元素模型描述单条肌肉–肌腱单元的力–长–速度关系。设第 i 块肌肉的最大等长收缩力为$F_{0i}$，标准化纤维长度$ \tilde{l}_i = l_i / l_{0,i}$，标准化收缩速度$\tilde{v}_i = v_i / v_{max,i}$，则其产生的肌腱拉力可表示为：
\begin{equation}
F_i^{M\pi} = F_{0j} \left[ a_i f_L(l_i) f_v(\tilde{v}_i) + f_P(l_i) + \beta \tilde{v}_i \right]
\label{eq:hill-muscle}
\end{equation}
式中 $aᵢ ∈ [0,1]$ 为肌肉激活水平；$fL(·)$、$fV(·)$、$fP(·) $分别为主动力–长度曲线、力–速度曲线与被动力–长度曲线；$β$为肌纤维阻尼系数。该模型同时刻画了主动收缩力、被动弹性力与速度阻尼三种力学贡献，能够较真实地反映人体上肢肌肉在驱动方向盘转动过程中的非线性力学特性。

神经控制信号 $uᵢ(t)$ 经过激活与去激活两段一阶动力学过程，方能形成实际的肌肉激活水平 $aᵢ(t)$，二者满足如下一阶非线性微分方程：
\begin{equation}
\dot{a}_i = \frac{u_i - a_i}{\tau(u_i, a_i)}
\label{eq:activation}
\end{equation}
其中时间常数 ${\tau(u_i, a_i)}$ 在激活相和去激活相分别取不同形式，以体现肌肉收缩–舒张的非对称特性：
\begin{equation}
\tau(u_i, a_i) =
\begin{cases}
\tau_{\text{act}} \left( 0.5 + 1.5 a_i \right) & , u_i \ge a_i \\[1em]
\dfrac{\tau_{\text{deact}}}{0.5 + 1.5 a_i} & , u_i \le a_i
\end{cases}
\label{eq:time-constant}
\end{equation}

本研究取激活时间常数 $\tau_{\text{act}}$ = 10 ms、去激活时间常数 $\tau_{\text{deact}}$ = 40 ms，与 MuJoCo 内置肌肉作动器默认值一致，可以较好地复现真实肌肉的响应延迟。激活动力学的存在意味着肌肉发力相对于神经指令存在固有的低通滤波效应，这是后续控制算法中需要考虑死区与饱和约束的生理依据之一。

在多刚体连杆模型下，全部肌肉对关节产生的合成力矩可由各自的肌腱拉力与瞬时力臂叠加得到：
\begin{equation}
\tau_{\text{mus}} = \sum_{i=1}^{m} r_i(q) \cdot F_i^{MT}
\label{eq:muscle-torque}
\end{equation}
式中$ m $为参与该关节运动的肌肉数量，$r_i(q)$ 为第 i 块肌肉关于广义坐标$ q $的瞬时力臂向量。该关系刻画了肌肉收缩力到关节驱动力矩的几何映射，是实现肌肉–方向盘耦合驱动的核心环节。
\section{MuJoCo 仿真引擎技术}
MuJoCo（ Multi-Joint Dynamics with Contact）是一款专门针对多关节接触系统设计的高性能物理仿真工具，由 Emo Todorov 研发而成，其核心应用领域集中在机器人与生物力学仿真，凭借高效的计算速度、精准的仿真效果以及强大的建模能力，成为该领域广泛应用的核心引擎之一。

该引擎的核心架构，采用广义坐标结合约束优化接触动力学的设计思路，有效弥补了传统物理仿真引擎的很多不足；在多体动力学求解方面，它以广义坐标为基础搭建系统模型，通过递归牛顿 - 欧拉算法高效求解运动方程，从根本上解决了笛卡尔坐标在数值计算时容易出现的不稳定问题；在接触求解环节，它不用传统的弹簧 - 阻尼仿真模型，而是通过凸优化方法精准处理接触约束，再结合层次包围盒与 GJK 算法实现快速碰撞检测，能够准确计算接触过程里的法向力和切向摩擦力，最大程度保证了仿真过程的真实与合理\cite{tu2026jiyu}。

在模型构建方面，MuJoCo 用 MJCF（XML）格式来描述模型，可以灵活支持骨骼、关节、肌肉、软组织等复杂生物力学结构的建模工作，同时也兼容 URDF 模型的直接导入，很适合人体上肢这类肌肉骨骼系统的仿真需求；引擎里自带专业的肌肉作动器模型，能实现肌肉激活信号到关节力矩的精准映射，为肌肉驱动和机械结构的耦合仿真提供了可靠的支撑\cite{liu2026noncontact}。

MuJoCo仿真引擎在数值计算层面具备明显的优化效果，该软件融入了稀疏矩阵优化、隐式积分等相关技术，能够兼容欧拉、龙格-库塔等多种常见的积分计算形式，这一特性可以有效保障整体的计算效率，也能大幅提高数值计算过程的稳定性；与此同时，该仿真引擎内置了正向动力学与逆向动力学两类计算功能模块，系统的运动状态可以依靠正向动力学完成求解计算，关节力矩数值则可以通过逆向动力学结合运动轨迹反向推导得出，这些基础功能都为本研究中肌肉驱动控制算法的设计与验证工作，提供了关键的技术支撑\cite{keaveny2026effects}。

MuJoCo的可视化效果依托OpenGL技术得以实现，该功能可以完成实时三维画面渲染和交互式调试工作，仿真过程中的肌肉激活状态、关节运动路径以及接触力的动态变化，都能够被直观展示出来，十分契合视觉引导操控的仿真试验需求；在本次整体研究工作中，高精度的仿真试验平台由MuJoCo仿真引擎搭建完成，平台实现了人体上肢肌肉骨骼模型与方向盘物理模型的组合建模，也能够为视觉引导操控试验、控制精度检测的全流程仿真验证工作，提供可靠的技术支撑。

在 MuJoCo 仿真引擎中，人体上肢被建模为多刚体系统，并以广义坐标 $q ∈ ℝⁿ $描述其位形。系统的前向动力学方程由经典拉格朗日形式给出：
\begin{equation}
M(q)\ddot{q} + C(q, \dot{q})\dot{q} + g(q) = \tau_{\text{mus}} + \tau_{\text{ext}}
\label{eq:lagrangian}
\end{equation}
式中 $M(q)$ 为关节空间惯性矩阵，$C(q, \dot{q})\dot{q}$ 为科里奥利力与离心力项，$g(q)$ 为重力项；$\tau_{\text{mus}}$ 由公式~\ref{eq:muscle-torque}给出，为肌肉合成的关节驱动力矩，$\tau_{\text{ext}}$ 为方向盘约束、外部接触等产生的外部力矩。MuJoCo 利用递归牛顿–欧拉算法高效求解 $M(q)$、$C(q, q̇)$、$g(q)$，避免了笛卡尔坐标下显式求解约束方程的数值不稳定问题。

为兼顾求解精度与计算效率，本系统采用 MuJoCo 默认的半隐式 Euler 数值积分方案推进仿真状态，其更新形式为：
\begin{equation}
q_{k+1} = q_k + \Delta t \cdot \dot{q}_{k+1}
\label{eq:euler}
\end{equation}
式中 $\Delta t$ 为仿真步长，本研究取 $\Delta t = 2 ms$（即 500 Hz 物理仿真频率）。半隐式 Euler 在保证较好稳定性的同时具有较低的计算开销，能够满足肌肉驱动方向盘操控仿真的实时性需求。
\section{User-in-the-Box 框架原理}
User-in-the-Box 是一种可以用来做人机交互任务仿真的开源框架，它采用模块化的设计方式，主要以生物力学用户模型作为核心内容，能够把人体肌肉骨骼模型、交互任务和感知模块快速组合到一起，形成可以单独运行的仿真环境，如图\ref{img-2}展示了这个框架里 Simulator 的组件关系，该框架把 MuJoCo 当作底层的仿真引擎，按照 OpenAI Gym 的接口标准进行设计，可以直接接入强化学习的训练过程，同时也支持把 OpenSim 模型转换成 MuJoCo 能够使用的肌肉骨骼模型，明显降低了生物力学建模的难度，从结构上来看，User-in-the-Box 包含三个关键的组成部分，分别是生物力学模型、交互任务和感知模块，生物力学模型主要用来完成人体运动的仿真，交互任务会对方向盘、操纵杆等操作对象进行定义，感知模块可以提供视觉、本体感觉等多种反馈信息，本研究直接按照这个框架的目录结构、配置文件要求和模型加载的逻辑，实现上肢肌肉模型和方向盘操控任务的快速整合，提高了系统开发的效率和规范程度。
\begin{figure}[htbp]
  \centering
  \includegraphics[width=0.8\textwidth]{image_3.png}
  \caption{Simulator的组件关系图}
  \label{img-2}
\end{figure}

User-in-the-Box 框架将一次完整的仿真单元抽象为生物力学模型$ B$、交互任务$ T $与感知模块$ P $三者的有序组合：
\begin{equation}
S = (B, T, P)
\label{eq:simulation}
\end{equation}
其中 $B$ 由 MuJoCo XML 文件描述上肢骨骼–肌肉系统，$T$ 描述方向盘等交互对象的几何与物理属性，$P$ 定义可获取的状态量及其编码方式。三者通过统一的 Simulator 类组装，对外暴露 OpenAI Gym 风格接口：
\begin{equation}
(o_{t+1}, r_t, d_t) = \text{step}(s_t, a_t)
\label{eq:gym-step}
\end{equation}
式中 $sₜ$ 为内部完整状态，$aₜ$ 为下达至各肌肉作动器的控制输入，$o_{t+1}$ 为感知模块返回的观测，$r_t$ 为任务回报，$d_t$ 为终止标志。对于本课题的方向盘操控任务，仅启用本体感觉与任务相关观测，因此观测向量定义为：
\begin{equation}
o_t = \left[ q_t, \dot{q}_t, a_t, x_t^{\text{obj}} \right]
\label{eq:observation}
\end{equation}
式中 $qₜ$、$\dot{q}_t$ 为关节角与角速度，$aₜ$ 为各肌肉当前激活水平，$x_t^{\text{obj}}$表示当前方向盘转角与目标转角。
\section{视觉引导与控制算法基础}
视觉引导可以提高人机交互的操控精准度和使用方便程度，是一项比较重要的技术，它能实时展示目标状态、当前状态和误差相关信息，帮助使用人员及时调整操作动作，从而实现稳定的跟踪效果，在方向盘操控系统当中，视觉引导一般包含动态目标角度生成、实时角度显示和误差可视化这几个部分，经常用到的目标轨迹有阶跃信号和正弦信号两种，这两种信号分别用来测试定点保持能力和连续跟踪能力，图~\ref{img-3}为视觉引导条件下肌肉驱动方向盘操控的闭环反馈结构，控制算法可以将肌肉激活信号转换成平稳、精准的设备运动状态，它的核心部分有信号滤波、非线性映射、增益调节和饱和约束，为了让操控的手感和稳定程度更好，系统一般会加入死区来过滤噪声，用低通滤波减少抖动，通过比例增益调整操作灵敏度，利用角度限幅避免出现超调的情况\cite{xie2026filament}。
\begin{figure}[htbp]
  \centering
  \includegraphics[width=0.8\textwidth]{image_4.png}
  \caption{视觉引导条件下肌肉驱动方向盘操控的闭环反馈结构图}
  \label{img-3}
\end{figure}

视觉引导以目标轨迹的形式向操作者实时呈现期望方向盘角度 $θ*(t)$。本研究采用阶跃信号与正弦信号两种典型轨迹，分别用于评估系统的定点保持能力与连续跟踪能力，其表达式如下：
\begin{equation}
\theta^{*}(t) = A \cdot 1(t \ge t_0)
\label{eq:step-trajectory}
\end{equation}

\begin{equation}
\theta^{*}(t) = A \cdot \sin(2\pi f t + \varphi)
\label{eq:sin-trajectory}
\end{equation}
式子中的 $A$ 代表目标幅值，$t_0$ 代表阶跃发生的时刻，$f$ 和 $\varphi$ 分别代表正弦轨迹的频率和初始相位，这两种轨迹会一起显示在仿真界面上，人们可以通过界面同时查看控制算法的瞬时反应情况和稳定后的跟踪效果，肌肉激活信号本身会出现高频抖动和基线漂移的情况，所以控制算法会先对原始信号 $u (k)$ 进行一阶指数加权移动平均的低通滤波处理，这种滤波方式也可以简称为 EMA：
\begin{equation}
u_f(k) = \alpha u(k) + (1-\alpha) u_f(k-1)
\label{eq:ema-filter}
\end{equation}
式中 $ \alpha \in (0, 1] $ 为平滑系数，$ \alpha $ 越小滤波越强但响应越慢。滤波后的信号继续经过死区处理以抑制零附近的随机抖动：
\begin{equation}
u_d = \text{sgn}(u_f) \cdot \max(|u_f| - u_{th}, 0)
\label{eq:deadzone}
\end{equation}
式中 $ u_{th} $ 为死区阈值，$ \text{sgn}(·) $ 为符号函数。死区之外，将处理后的信号经非线性增益函数 $ g(·) $ 与比例增益 $ K_p $ 映射至方向盘转角指令，并由饱和函数限制在物理可达范围之内：
\begin{equation}
\theta_{\text{cmd}} = \text{sat}\left( K_p \cdot g(u_d), -\theta_{\max}, \theta_{\max} \right)
\label{eq:command}
\end{equation}

其中饱和函数定义为：
\begin{equation}
\text{sat}(x, l, h) = \max\left( \min(x, h), l \right)
\label{eq:saturation}
\end{equation}
$\theta_{\max}$ 指的是方向盘可以转动的最大角度，这种映射方式可以同时照顾到操作灵敏度和使用安全，在发力较小的区域，非线性增益会把细微的动作放大，以此提高小角度操控的准确程度，在发力较大的区域，饱和约束可以避免出现超调和设备损坏的情况，为了让方向盘可以按照闭环方式跟踪目标轨迹，需要对瞬时跟踪误差进行定义：
\begin{equation}
e(t) = \theta^{*}(t) - \theta(t)
\label{eq:tracking-error}
\end{equation}

该误差量同时作为视觉引导界面的实时反馈量与第 2.5 节中精度指标的计算输入，是连接控制算法与系统评估的关键桥梁。
\section{系统评估指标体系}
为全面、客观、量化地评价肌肉驱动生物力学人机交互系统的控制效果与运行性能，本研究构建多维度、模块化的评估指标体系，从数据质量、运动学解算、控制精度、仿真稳定性、综合任务效果五个维度开展评估。
\begin{table}[htbp]
  \centering
  \caption{系统评估指标体系}
  \label{tab:eval_index}
  \begin{tabular}{p{2.5cm}p{3.5cm}p{8cm}}
    \toprule
   模块 & 评估维度 & 核心指标 \\
    \midrule
    数据处理模块 & 数据质量 & 数据同步误差、预处理延迟、数据缺失率 \\
    运动学模块 & 解算精度与合理性 & 关节位置误差、姿态平滑度、关节超限率 \\
    控制模块 & 跟踪与控制精度 & 平均绝对误差 (MAE)、均方根误差 (RMSE)、响应延迟、最大误差 \\
    仿真模块 & 运行稳定性 & 仿真帧率稳定性、无故障运行时长、角度漂移量 \\
    评估模块 & 综合任务效果 & 轨迹相似度、生物合理性、目标跟踪成功率、任务完成度 \\
    \bottomrule
  \end{tabular}
\end{table}

核心指标说明如下：

平均绝对误差 MAE 可以用来表示方向盘角度跟踪的整体偏差情况，数值越小就代表控制的精准程度越高，均方根误差 RMSE 对于比较大的误差会更加敏感，人们常用它来判断系统控制的稳定程度和波动情况，响应延迟是从肌肉激活信号输入，到方向盘做出相应转动输出的时间间隔，能够体现系统的实时运行能力，角度漂移量是系统在长时间运行时，方向盘在没有控制输入的情况下出现的角度偏移数值，可用于评价系统的稳态稳定性能，轨迹相似度用来对比实际运动轨迹和目标运动轨迹的接近程度，能够体现视觉引导条件下的跟踪效果，生物合理性用于判断运动过程是否符合人体上肢肌肉和关节的正常运动规律，其中也包括关节是否超出限制、肌肉发力是否自然等内容，任务完成度是在规定时间内达到目标角度、保持稳定姿态并且完成操控任务的成功概率，为了能够定量描述系统的方向盘控制精准程度，我们将采样总数设为 N，第 i 个采样点对应的目标角度与实际角度分别设为 $\theta_i^*$ 和 $\theta_i$，平均绝对误差MAE可以按照如下方式进行定义：
\begin{equation}
\text{MAE} = \frac{1}{N} \sum_{i=1}^{N} |\theta_i^{*} - \theta_i|
\label{eq:mae}
\end{equation}
MAE反映方向盘角度跟踪的整体偏差水平，其值越小表示控制精度越高。在此基础上，对较大误差更加敏感的均方根误差（RMSE）定义为：
\begin{equation}
\text{RMSE} = \sqrt{\frac{1}{N} \sum_{i=1}^{N} \left( \theta_i^{*} - \theta_i \right)^2}
\label{eq:rmse}
\end{equation}
RMSE在统计意义上等价于误差的二阶矩，能够有效反映系统控制的稳定性与波动程度。为进一步刻画极端误差，定义最大误差为：
\begin{equation}
e_{\max} = \max_i \left( \left| \theta_i^{*} - \theta_i \right| \right)
\label{eq:max-error}
\end{equation}
$e_{\max}$ 用于判定系统在跟踪过程中是否出现单点超限现象，是评估安全性的重要指标。

对于响应延迟，本研究采用实际角度序列 $θ(t)$ 与目标角度序列 $θ^*(t)$ 之间互相关函数取最大值时对应的时间偏移量作为定义：
\begin{equation}
t_d = \arg\max_{\tau} \left( \sum_{t} \theta(t+\tau) \cdot \theta^{*}(t) \right)
\label{eq:delay}
\end{equation}
式中 $t_d$ 即为系统从肌肉激活信号输入到方向盘输出对应转角所需的等效延迟时间。综合上述指标，可从均值偏差、波动程度、极端误差与时滞特性四个维度对肌肉驱动方向盘操控系统做出量化评估。
